% The comparison table, kept in its own file so it can be regenerated or rewritten without
% touching the prose around it. Included by comparison.tex.
%
% One row per method (or per group of methods that agree in every descriptive column).
% Every row needs a reference: never add an entry whose source you cannot cite. Rename,
% add or drop columns to fit the topic, and keep the tabularx column spec in step.
% [!t]: a two-column-wide float can only sit at the top of a page. Keeping the \input at the start
% of the section, before the figure and the family subsections, makes it land on the first page of
% the section instead of drifting to the end of the paper.
\begin{table*}[!t]
\caption{Summary of the methods. \todo{Describe what a row is and what the columns mean.}}
\label{tab:comparison}
\small
\begin{tabularx}{\textwidth}{@{}L L L L l@{}}
\toprule
\textbf{Input / data used} & \textbf{Output / quantity} & \textbf{Model or assumption} &
\textbf{Technique} & \textbf{Ref.}\\
\midrule
\todo{input} & \todo{output} & \todo{model} & \todo{technique} & \cite{key1}\\
\addlinespace
\todo{input} & \todo{output} & \todo{model} & \todo{technique} & \todo{ref}\\
\bottomrule
\end{tabularx}
\end{table*}
